%% 301_Information_Umkehrtest_De_ch.tex
%% FFGFT -- Dok. 301: Der Umkehrtest -- Information als Ausgang, nicht als Grundeinheit;
%%                    das Bit als ultimative, eingangsbedingungen-tilgende Reduktion
%% Autor: Johann Pascher
%% Datum: Juli 2026
%% Sprache: Deutsch

\section*{Dok.~301 \quad Der Umkehrtest --- Information als Ausgang, nicht
als Grundeinheit}

\section{Anlass und Abgrenzung}

Aus der Zusammenarbeit stammt ein Stabilitätsbild, das drei Größen ---
Information, Kohärenz und Energie (ICE) --- als gleichrangige Ecken eines
Dreiecks führt, mit einer Referenz im Zentrum. Dieses Dokument prüft, welche
kategoriale Stellung \enquote{Information} in einem solchen Bild tatsächlich
einnimmt. Der Befund ist gerichtet, aber nicht einstufig: \emph{Information}
ist nicht gleichbedeutend mit \emph{Bit}. Der Begriff zerfällt in drei Ebenen ---
das Bit (Comparator-Ausgang), die Beschreibung (gewählte Darstellung) und die
absolute Identität (das vollständige Objekt selbst). Als Grundeinheit neben
Energie taugt Information auf keiner der beiden unteren Ebenen; fundamental gilt
sie allein im absoluten Sinn, als oberste Identität. Das Kriterium, das diese
Einordnung trägt, ist die Umkehrbarkeit. Abgrenzung: Kategorie-B-Befund
(kategoriale Klassifikation, keine physikalische Ableitung, keine Behauptung,
dass FFGFT ein Bit herleitet). Das Ergebnis ist eine Präzisierung des
Reduktionsprinzips aus Dok.~296 und der Comparator-Fassung aus Dok.~300.

\section{Das Kriterium: fundamental heißt umkehrbar}

Eine Größe ist genau dann fundamental, wenn ihre Umformung keine
ausgezeichnete Richtung hat --- wenn man sie auf den Kopf stellen kann und sie
stehen bleibt. Der Prüfstein ist die Energie in natürlichen Einheiten. Über
die Relation
\begin{equation}
  \tilde{T}\cdot m = 1
\end{equation}
sind Energie, Masse und die duale Zeit keine unabhängigen Achsen, sondern drei
Ansichten einer Größe:
\begin{equation}
  E \;=\; m \;=\; \frac{1}{\tilde{T}}.
\end{equation}
Diese Umformung ist symmetrisch: man kann bei $E$ beginnen und $m$ ablesen,
bei $m$ und $\tilde{T}$, bei $\tilde{T}$ und $E$. Keine Seite ist die primäre,
keine ist Eingang oder Ausgang. Genau das ist der Grund, warum Energie
intrinsisch ist und die Energieskala $E_0$ dem Rahmen nativ zukommt (Dok.~292):
die Umformung führt nichts weg, jede Ansicht trägt die anderen vollständig in
sich. Umkehrbarkeit ist hier kein Zusatz, sondern Ausdruck davon, dass nichts
verloren geht. In SI treten $\hbar$ und $c$ als Umrechnungsfaktoren wieder ein
und ziehen die drei Größen auseinander; das ist die Projektion, nicht der Kern.

\section{Intrinsisch und relational}

Nicht jede Größe trägt sich selbst. Eine geometrische Größe ist
\textbf{intrinsisch}: eine Fläche, eine Windung hat ihren Wert ohne Bezug auf
etwas anderes. Eine thermodynamische oder informationelle Größe ist
\textbf{relational}: Entropie ist nur relativ zu einer Vergröberung
(Coarse-Graining) definiert, ein Bit nur relativ zu einer ausgezeichneten
Referenz. Es gibt kein Bit \enquote{an sich}. Diese Unterscheidung ist die
kategoriale Grundlage des Folgenden: intrinsische Größen können fundamental
sein, relationale setzen stets ein \emph{bezüglich} voraus und können daher
nicht selbsttragend sein.

\section{Das ICE-Dreieck als Comparator}

Liest man das ICE-Bild schaltungstechnisch, deckt es sich Ecke für Ecke mit
dem Comparator aus Dok.~300 --- nicht analog, sondern als dasselbe Objekt von
der anderen Seite. Verwertbare Information am Ausgang eines Vergleichs verlangt
drei Zutaten und eine vierte Bedingung:
\begin{itemize}
  \item \textbf{Energie} --- das Signal am $+$Eingang trägt einen Pegel;
    ohne Pegel kein Vergleich.
  \item \textbf{Kohärenz} --- das Signal hat Phasenstruktur, ist nicht bloßes
    Rauschen; geometrisch die Phase auf den Windungen.
  \item \textbf{Referenz} --- der $-$Eingang, die Schwelle, gegen die
    entschieden wird; er muss aus einer \emph{anderen}, abgeleiteten Quelle
    stammen als das Signal (Dok.~300: Referenz aus Ableitung, nicht aus Kopie).
  \item \textbf{Hysterese} --- eine nackte Schwelle klappert am Umschaltpunkt;
    ein Übergang gilt erst, wenn er über ein Totband steht (Amplituden-Hysterese)
    oder über ein Fenster stehen bleibt (Zeit-Hysterese, Entprellung).
    Ohne Totband ist der Ausgang bedeutungslos, gleich wie sauber Energie und
    Kohärenz sind.
\end{itemize}
\emph{Wichtig:} Das ICE-Dreieck enthält Energie, Kohärenz und Referenz, aber
kein ausgewiesenes Totband. Die Hysterese ist die vierte Bedingung, ohne die
aus \enquote{Schwelle überschritten} kein \enquote{Übergang gilt} wird. Die in
der finiten Proxy-Implementierung verwendete Zeit-Hysterese ist genau das
gleitende Fenster $M_m$ (Dok.~300/297) --- ein Entpreller, und als solcher ein
Kurzzeitpuffer, kein akkumulierendes Gedächtnis.

\section{Information ist Ausgang, nicht Eingang}

Damit dreht der Comparator die Ordnung um, die das Dreieck nahelegt. Energie,
Kohärenz, Referenz und Hysterese stehen am \emph{Eingang}. Information ist das,
was am \emph{Ausgang} herauskommt, wenn diese Bedingungen erfüllt sind: das
Bit, das der Schwellenentscheid erzeugt. Sie ist nicht die Zutat, aus der der
Vergleich gebaut ist, sondern sein Ergebnis. Eine Größe aber, die den ganzen
Apparat voraussetzt, um zu entstehen, kann nicht die Grundeinheit sein, aus der
der Apparat besteht. Das ist zugleich der Grund, warum frühere Versuche, ein
Bit als Grundeinheit zu verankern, scheitern mussten: sie wollten das
Ausgangsprodukt als Eingangsgröße setzen --- eine Zirkularität.

\section{Der Umkehrtest}

Das Kriterium aus Abschnitt~2, auf Information angewandt, fällt negativ aus.
$E$--$m$--$\tilde{T}$ ist bidirektional: jede Ansicht rekonstruiert die
anderen, weil nichts weggeworfen wird. Der Comparator ist \emph{unidirektional}:
\begin{equation}
  \text{Signal} + \text{Referenz} + \text{Hysterese}
  \;\longrightarrow\; \text{Bit},
\end{equation}
und dieser Pfeil kehrt nicht um. Aus einem Bit erhält man das Signal nicht
zurück, nicht seinen Pegel, nicht seine Phase, nicht die Referenzschwelle, gegen
die entschieden wurde. Fundamentale Größen sind umkehrbar, weil sie erhalten;
Information ist nicht umkehrbar, weil sie tilgt. \emph{Fundamental und Bit
schließen sich damit nicht zufällig aus, sondern definitorisch}: Umkehrbarkeit
verlangt, dass die erzeugende Struktur im Ergebnis erhalten bleibt, und das Bit
ist gerade die Operation, die sie nicht erhält.

\section{Das Bit als ultimative Reduktion}

Nach Dok.~296 ist jede Auswertung eine Reduktion: sie lässt weg, und gerade das
Weglassen macht einen Teilbereich verwertbar. Das Bit ist der Grenzfall dieser
Aussage --- die \emph{maximale} Reduktion. Ein ganzer Signalverlauf mit seiner
Energie, seiner Phase, seiner Lage relativ zur Schwelle wird auf ein einziges
Ja/Nein zusammengezogen; weniger geht nicht. Und die Irreversibilität ist nicht
ein Zusatz zur Reduktion, sie \emph{ist} die Reduktion: die Schwelle, die Phase,
die Eingangsbedingungen \emph{waren} --- im Moment, in dem der Comparator kippt,
sind sie weg. Das Bit trägt nicht mehr, wie es zustande kam. Es sagt
\enquote{drüber} oder \enquote{drunter}, nicht mit welchem Pegel, nicht mit
welcher Phase, nicht wie weit von der Referenz. Alles, was das Ergebnis erzeugt
hat, ist im Ergebnis nicht mehr enthalten --- deshalb gibt es kein Zurück, denn
das Wozu-zurück ist genau das, was die Reduktion getilgt hat.

\section{Der Landauer-Anschluss}

Diese Irreversibilität ist die thermodynamische Seite derselben Sache. Ein Bit
zu löschen kostet $k_B T \ln 2$ und erzeugt Entropie; die Richtung ist
eingebaut, weil das Tilgen von Information Wärme produziert. Den Comparator
umzukehren hieße, diese Entropie zurückzuholen, und das ist thermodynamisch
ausgeschlossen. Das erklärt rückwirkend, warum die Anbindung eines
Grundeinheit-Bits an die Wärmebewegung im bisherigen Arbeitsstand so schlecht
gelang: der Versuch band eine gerichtete, reduzierende Größe an eine reversible
Grundstruktur; die Richtungen passen nicht zusammen. Nicht ein Mangel an
Rechnung, sondern eine kategoriale Unverträglichkeit war die Ursache.

\section{Drei Ebenen von Information: Bit, Beschreibung, absolute Identität}

Der Umkehrtest trennt sauber, aber die Grenze verläuft nicht bei
\enquote{Information} pauschal. Der Begriff zerfällt in drei Ebenen, die
verschieden auf den Test antworten.

\emph{Erste Ebene --- das Bit.} Der Comparator-Ausgang: referenzabhängig, am
Ausgang erzeugt, eingangsbedingungen-tilgend. Er fällt durch --- unidirektional,
reduzierend, die ultimative Reduktion (siehe \enquote{Der Umkehrtest} und
\enquote{Das Bit als ultimative Reduktion}). Folge, nicht Grundlage.

\emph{Zweite Ebene --- die Beschreibung.} Information ist nicht das Bit, sondern
zunächst eine \emph{Beschreibung}: die Darstellung eines Sachverhalts. Welche
Art der Beschreibung gewählt wird --- geometrisch, mathematisch, rhetorisch ---
ändert die Darstellung, nicht den Sachverhalt. Das ist dieselbe Linie wie in
Dok.~297/300: eine Darstellung ist eine Koordinatenwahl, keine strukturelle
Eigenschaft des Objekts; verschiedene Beschreibungen desselben Objekts sind
ineinander übersetzbar, aber keine ist das Objekt. Und keine ist je
\emph{vollständig}: jede Beschreibung ist eine Projektion (Dok.~296), also eine
Reduktion, also lässt sie weg. Eine vollständige Beschreibung gäbe es nur, wenn
eine Projektion nichts wegließe --- und dann wertete sie nichts aus. Auch die
Beschreibung fällt daher durch den Umkehrtest: als Weg zurück zum Objekt ist
keine einzelne Beschreibung umkehrbar, weil jede unvollständig ist. Die
Beschreibungen sind untereinander übersetzbar, aber keine rekonstruiert das
Ganze.

\emph{Dritte Ebene --- die absolute Identität.} Fundamental im vollen Sinn ist
allein das vollständige Objekt selbst: das randlose $T^4$-Objekt (der Kern
$\xi$ auf $T^4$), aus dem jede Projektion nur einen Ausschnitt zeigt.
\emph{Information im absoluten Sinn} ist keine Größe \emph{im} System und keine
einzelne Beschreibung, sondern die vollständige Selbstidentität des Objekts mit
sich, vor jeder Projektion. Nur diese Ebene besteht den Umkehrtest --- trivial,
weil bei ihr nichts wegzulassen ist: es gibt kein \enquote{bezüglich}, weil es
keinen Außenstandpunkt gibt. Das ist die oberste Identität; ihr, und nur ihr,
kommt \enquote{Information} fundamental zu (Selbstreferenz, Dok.~248 §8).

\emph{Wichtig:} Die geometrische Struktur des Kerns --- Windungszahlen als
topologische Invarianten, verlustfrei getragen --- gehört zur dritten Ebene, zum
Objekt selbst; sie ist nicht \enquote{Information} im Bit-Sinn. Eine geometrische
\emph{Beschreibung} dieser Struktur bleibt dagegen zweite Ebene: eine Projektion.
Und das abgelesene Bit ist erste Ebene. Die drei Bedeutungen dürfen nicht
vermischt werden. Insbesondere gilt: steht überhaupt etwas \enquote{über} der
Energie, dann nicht das Bit, sondern das vollständige Objekt --- und Energie ist
eine seiner Projektionen (eine besonders intrinsische, aber eine Ansicht).

\section{Digitalisierung: die Stufung und ihre Glättung}

Der Grund, warum ein Bit nie die gesamte Information tragen kann, ist der
Grundmechanismus der Digitaltechnik: Analog-Digital- und Digital-Analog-Wandlung.
Ein Bit ist die 1-Bit-Quantisierung --- die gröbste mögliche Stufung, eine
einzige Schwelle. Jede feinere Stufung mit endlich vielen Stufen bleibt
kategorial dasselbe Objekt, nur höher aufgelöst.

Die Torsionswindung liefert das natürliche Beispiel. Die aufgerollte Rekursion
schließt nach genau $75$ gleichen Teilstücken (Dok.~300/295); zerlegt man die
Windung in diese $75$ Teile, hat man sie digitalisiert --- $75$ diskrete Stufen.
Per DA-Wandlung lässt sich daraus wieder ein kontinuierlicher Verlauf gewinnen.
Aber die Rekonstruktion ist bereits verlustbehaftet, aus zwei Gründen:
\begin{itemize}
  \item Die Quantisierung wirft weg, was \emph{zwischen} den Stufen lag. Die
    kontinuierliche Struktur zwischen den $75$ Knoten --- die exponentiell
    schrumpfenden Schritte der log-Spirale, das kumulierte Integral $D(k)$
    (Dok.~295) --- ist in den $75$ Werten nicht enthalten.
  \item Der DA-Wandler gibt eine Treppe aus, keine glatte Kurve. Man muss die
    Stufen \emph{glätten} (Rekonstruktionsfilter). Aber die Glättung ist eine
    Interpolation --- eine Annahme, die der Filter hinzufügt, nicht eine im
    Digitalwert enthaltene Information. Was zwischen den Stufen steht, wird
    geraten, nicht zurückgewonnen.
\end{itemize}
Damit ist der AD/DA-Umlauf gerade \emph{kein} echter Kehrwert, sondern ein
gerichteter, verlustbehafteter Vorgang: der Umkehrtest in schaltungstechnischer
Form. Die verlorene Zwischenstruktur wird durch die Glättung nicht
rekonstruiert, sondern durch eine Filterannahme ersetzt --- dieselbe Signatur wie
beim Bit, dessen Eingangsbedingungen \enquote{waren} und im Ergebnis nicht mehr
stehen.

\emph{Wichtig:} Selbst die \emph{natürliche} Stufung --- $75$, die
Schließungszahl --- trägt das Kontinuum nicht. Sie ist die richtige Zahl für den
Schließungsschritt, aber nicht für den kontinuierlichen Verlauf: das volle
Objekt ist reicher als jede endliche Stufung (Dok.~296). Kein $N$, auch nicht das
strukturell ausgezeichnete $N=75$, trägt das Kontinuum. Eingeordnet in die drei
Ebenen: das Bit ist die 1-Stufen-Quantisierung (erste Ebene, Ausgang); jede
endliche $N$-Stufung --- auch die $75$-Teilung --- ist eine Beschreibung endlicher
Auflösung (zweite Ebene), feiner als ein Bit, aber weiter verlustbehaftet, weiter
Projektion; allein die kontinuierliche Windung, das vollständige Objekt, trägt
alles (dritte Ebene). Die Glättung ist genau die Stelle, an der die Beschreibung
ihre Unvollständigkeit durch eine Annahme überdeckt.

\section{Präzisierung: Resonanzen erzwingen die Quantisierung}

Der vorige Abschnitt sprach von der \enquote{kontinuierlichen Struktur zwischen
den Knoten}. Das ist zu präzisieren, denn auf der feinsten Skala ist das Objekt
nicht kontinuierlich. Der Torus ist kompakt, und ein kompakter Torus hat ein
\emph{diskretes} Spektrum: die Resonanzen auf $T^4$ --- die Eigenmoden des
$T^4$-Connection-Laplace (Dok.~283) --- erzwingen eine Quantisierung. In den
kleinsten Teilen sind \emph{nicht alle Möglichkeiten gleichberechtigt}; die
Resonanzstruktur zeichnet bestimmte diskrete Zustände aus. Genau das ist der
Grund, warum das Standardmodell diskret herauskommt --- die Massen, $\theta=2/9$,
die $Z_3$-Struktur sind Resonanzen, keine beliebigen Punkte eines Kontinuums.

Daraus folgt eine Einschränkung des Verlust-Arguments: es gibt einen
\emph{kleinstmöglichen Sprung}, das von den Resonanzen gesetzte Quantum. Eine
Wandlung, die bis zu diesem kleinsten Sprung geht und die Resonanzstruktur
trifft, ist \emph{erlaubt} --- sie ist nicht verlustbehaftet, weil unterhalb des
Sprungs keine weiteren Zustände liegen, die man verlieren könnte. Was zwischen
einer \emph{gröberen} Stufung liegt, ist daher kein Kontinuum, sondern feinere
\emph{diskrete} Resonanzstruktur, und diese hat einen Boden.

Damit verfeinert sich der Umkehrtest: Verlust entsteht nicht durch Quantisierung
als solche, sondern nur durch eine Stufung, die \emph{gröber} als die Resonanzen
ist oder \emph{quer} zu ihnen liegt --- etwa eine beliebige Gleichteilung. Eine
resonanzgetreue Wandlung bis zum kleinsten Sprung besteht den Umkehrtest: sie
erfasst die tatsächliche diskrete Struktur, nicht eine Näherung. Das Bit bleibt
gleichwohl die ultimative Reduktion --- eine einzige Schwelle ist gröber als der
Resonanzboden und trifft ihn nicht. Und über eine einzelne Skala hinaus gilt
weiterhin: keine Beschreibung einer Skala erfasst die skalenübergreifende Fülle
des vollen Objekts (Dok.~296). Die Aussage \enquote{das volle Objekt ist reicher
als jede endliche Stufung} ist damit \emph{skalenübergreifend} zu lesen ---
innerhalb einer Skala ist die resonanzgetreue Stufung dagegen die wahre
Struktur.

\section{Zeitliche und topologische Randlosigkeit}

Zwei Eigenschaften der absoluten Identität folgen aus Analysis und Topologie und
sind daher als \emph{Folgerungen} gebucht, nicht als Setzung.

\emph{Kein erreichbarer Rand der Rekursion.} Die Abbildung
$\xi_{n+1}=\xi_n(1-100\,\xi_n)$ hat $\xi=0$ als Fixpunkt; $\xi_n$ zerfällt monoton
gegen null, erreicht null aber in keinem endlichen Schritt (Doc.~296). Es gibt
keinen ersten und keinen letzten Schritt. Ein \enquote{einfachster Anfangszustand}
--- eine Skala ohne fraktale Struktur, aus der die Verfeinerung erst zu wachsen
begänne --- existiert nicht: die Selbstähnlichkeit (Faktor $1/e$ je $2\pi$-Umlauf,
log-Spirale) ist keine Eigenschaft, die sich aufbaut, sondern in der
Abbildungsvorschrift bereits vollständig vorhanden. Ein struktureller Anfang würde
die Skaleninvarianz brechen, die die Rekursion definiert --- ein Widerspruch, kein
Weltbild.

\emph{Zeit ohne markierbaren Ursprung.} Die erlebbare Zeit \emph{ist} der
kumulierte Defekt $D(k)=\sum 100\,\xi_n$ (Doc.~295), das Integral über
$\delta t = 1/(100\,\xi_n)$. Da die Rekursion keinen ersten Schritt hat, hat
$D(k)$ keinen ausgezeichneten Ursprung. Zeit ist die Rekursion, nicht ihr
Behälter; einen \enquote{Anfang der Zeit} zu suchen hieße, die Rekursion an einen
Zeitpunkt zu setzen, den sie erst erzeugt.

\emph{Kein Außenstandpunkt.} $T^4$ ist kompakt und randlos, $\partial T^4 =
\emptyset$. Das ist eine topologische Tatsache, keine Setzung. Aus ihr folgt, dass
es keinen Standpunkt außerhalb des Objekts gibt --- und damit das \enquote{kein
bezüglich} der absoluten Identität: das vollständige Objekt hat kein Außen, gegen
das es sich abgrenzte.

Beides zusammengenommen: die absolute Identität ist randlos im Räumlichen (kein
Außen, $\partial T^4=\emptyset$) und im Zeitlichen (kein Anfang, keine erreichbare
Grenze der Rekursion) --- dieselbe Eigenschaft in zwei Gestalten: kein Standpunkt
jenseits des Objekts. Diese Aussagen sind aus Analysis und Topologie ableitbar und
tragen keinen Spekulationsmarker.

\section{Der reflexive Fall --- was über die Mathematik hinausgeht}

\emph{Das Folgende geht ausdrücklich über die Mathematik hinaus und ist als
philosophisch-spekulativ gebucht (P35). Es sind zwei Interpretationssprünge, keine
Rechnungen.}

Wendet man die drei Ebenen auf FFGFT selbst an, so ist das Framework --- der Torus,
$\tilde{T}\cdot m=1$, $\xi$ über alle Skalen --- eine \emph{Beschreibung} (zweite
Ebene), notwendig unvollständig. Das ist keine Schwäche gerade dieses Frameworks,
sondern was eine Beschreibung ist: keine Theorie ist je ihr Objekt, jede projiziert
und lässt weg. FFGFT sagt dies ausdrücklich, während naive Realismen so tun, als
seien sie das Gebiet.

\emph{Erster Sprung --- die Realisten-Setzung.} Dass die resonanzgetreue
Beschreibung ein \emph{reales} Objekt nachspurt und nicht bloß eine konsistente
Struktur ist, ist gesetzt, nicht bewiesen. Die Mathematik zeigt Resonanztreue
($\alpha$ beim Messwert, die Massen als Eigenmoden, Koide); dass Treue Realität
nachspurt statt nur Konsistenz, ist die Realisten-Wette. Der deflationäre Leser darf
stets sagen: alles nur Beschreibung, das Objekt eine nützliche Fiktion. FFGFTs
Antwort ist kein Beweis, sondern die Wette, dass der resonanzgetreue Erfolg durch
ein reales Objekt besser erklärt ist als durch Zufall --- und das Besondere ist nur,
dass die Resonanztreffer diese Wette falsifizierbar machen.

\emph{Zweiter Sprung --- die Diagnose des Wort-Reflexes.} Die Anziehungskraft der
Information-als-Fundament-Idee (\enquote{Im Anfang war das Wort}; \enquote{it from
bit}; die Masse-Energie-Information-Äquivalenz) könnte aus einem tiefen kulturellen
Reflex stammen: dass Information die erste Substanz sei, aus der Materie folgt. Der
Umkehrtest dreht die Reihenfolge um: das Wort ist das Bit, und das Bit ist die
letzte Reduktion, nicht die erste Schöpfung. Im Anfang war die Geometrie (das
Objekt); das Wort ist, was eine ihrer Skalen davon ausliest. Dies ist eine
kulturell-kognitive Hypothese darüber, \emph{warum} die Intuition zieht --- keine
mathematische Aussage; der Umkehrtest selbst bleibt mathematisch.

\emph{Selbstprojektion.} FFGFT ist ein $T^4$-Prozess, der $T^4$ beschreibt
(Doc.~248 §8): die Beschreibung ist aus demselben Stoff wie das Objekt und wird von
ihm erzeugt (der beschreibende Prozess ist selbst ein $T^4$-Prozess auf der
kognitiven Skala). Damit ist FFGFT weder extern zum Objekt noch das Objekt, sondern
dessen \emph{Selbstprojektion} --- das Objekt, das sich auf einer seiner eigenen
Skalen beschreibt. Selbstreferenz als ontologischer Ort, gebucht als Setzung.

Daraus die Grenze der Ausgangsfrage: \enquote{Wieviel ontologische Realität hat
FFGFT?} ist selbst eine Frage der zweiten Ebene --- sie setzt einen beschreibenden
Standpunkt voraus, den es gegenüber der dritten Ebene nicht gibt. Man kann nicht von
außen fragen, wieviel Realität das Objekt hat, weil es kein Außen gibt; man kann
nicht fragen, was vor ihm war, weil es kein Vorher gibt. Auf der Ebene der absoluten
Identität löst die Frage sich auf.

\section{Status}

Kategoriale Klassifikation (Kategorie~B), kein physikalischer Ableitungsschritt.
Der Befund ordnet Information relativ zu Programmen ein, die sie als
physikalische Grundgröße setzen (etwa eine Masse-Energie-Information-Äquivalenz):
solche Programme behandeln Information als Eingang; der Umkehrtest zeigt, dass
das Bit Ausgang ist. Die drei Ebenen sind dabei mit verschiedenen Temperaturen
zu buchen:
\begin{itemize}
  \item \textbf{Bit} --- \emph{bewiesen}: die ultimative, eingangsbedingungen-tilgende
    Reduktion ist irreversibel und daher Folge, nicht Grundeinheit (Umkehrtest).
  \item \textbf{Beschreibung} --- \emph{strukturell}: jede Darstellung ist eine
    Projektion und damit unvollständig (Reduktionsprinzip, Dok.~296); eine
    Koordinatenwahl, keine strukturelle Eigenschaft des Objekts.
  \item \textbf{Absolute Identität} --- \emph{ontologische Setzung}: Information
    im absoluten Sinn wird mit dem vollständigen Objekt identifiziert und ist die
    oberste Identität (Selbstreferenz, Dok.~248 §8). Das ist gesetzt, nicht
    abgeleitet.
\end{itemize}
Die tragende Aussage in einem Satz: \emph{fundamental $\Leftrightarrow$
umkehrbar; Bit und einzelne Beschreibung fallen durch (das Bit tilgt, jede
Beschreibung ist unvollständig), allein die absolute Identität besteht den Test
--- also gilt Information nur im absoluten Sinn als oberste Identität, und diese
ist bei FFGFT der Kern $\xi$ auf $T^4$, nicht eine Größe daneben}. Anschluss:
Dok.~296 (Reduktion als Bedingung der Auswertbarkeit), Dok.~300 (Comparator,
Referenz, Hysterese), Dok.~283 (diskretes Resonanzspektrum auf $T^4$),
Dok.~292 ($E_0$ intrinsisch), Dok.~248 §8 (Selbstreferenz).
